\section{DISEÑO DISCRETO DE CASOS Y PLAN DE COBERTURA}
\label{anx:diseno-casos}
\thispagestyle{empty}

La Tabla~\ref{tab:diseno-casos} presenta el diseño discreto de casos y los parámetros físicos de cobertura.

Este anexo define cómo se formarán VI3 y VI4 y cómo se cubrirá el espacio físico de los casos. VI3 combina el artefacto verificado VD2 con el catálogo de escenarios $E_3$; VI4 combina el expediente térmico VD3 con el catálogo de representaciones $C_4$. El contraste, la extensión y la proximidad son parámetros que caracterizan un escenario ya seleccionado; por ello, sus niveles forman una malla de sensibilidad física y no un diseño factorial de variables independientes.

\scriptsize
\setlength{\tabcolsep}{3pt}
\renewcommand{\arraystretch}{1.18}
\begin{longtable}{L{3.10cm}L{1.45cm}L{3.00cm}L{5.15cm}L{5.25cm}}
\caption{Diseño discreto de casos y parámetros físicos de cobertura. Fuente: Elaboración propia.}
\label{tab:diseno-casos}\\
\toprule
\textbf{Elemento} &
\textbf{Código} &
\textbf{Papel metodológico} &
\textbf{Categorías o niveles iniciales} &
\textbf{Uso en el plan de cobertura} \\
\midrule
\endfirsthead

\toprule
\textbf{Elemento} &
\textbf{Código} &
\textbf{Papel metodológico} &
\textbf{Categorías o niveles iniciales} &
\textbf{Uso en el plan de cobertura} \\
\midrule
\endhead

\midrule
\multicolumn{5}{r}{\small Continúa en la siguiente página.}\\
\midrule
\endfoot

\bottomrule
\endlastfoot

Artefacto verificado heredado & VD2 / $\mathrm{VI}_{3,h}$ & Producto de OE2 y componente obligatorio de VI3 & Aceptado en el dominio; aceptado con restricciones compatibles. El estado rechazado no pasa a OE3. & Fija la versión del artefacto y su dominio de validez antes de recorrer escenarios. \\

Catálogo nuevo de escenarios & $E_3$ & Componente nuevo de VI3 & Base homogéneo; estratificado; inclusión localizada; material de mejora; zona seca; combinación multimaterial. & Selecciona el objeto físico-computacional; cada realización contiene mapa $\kx$ y ficha de parámetros. \\

Expediente de evaluación térmica & VI3 & Variable independiente discreta de OE3: $\{\mathrm{VD2},E_3\}$ & Combinaciones válidas de una VD2 aceptada con una realización de $E_3$. & Constituye la entrada completa de OE3 y produce VD3. \\

Expediente térmico heredado & VD3 / $\mathrm{VI}_{4,h}$ & Producto de OE3 y componente obligatorio de VI4 & Sin efecto relevante; efecto relevante no crítico; condición crítica. El estado inconcluso no pasa a OE4. & Conserva mapas, métricas, límites y clasificación del escenario que será comparado. \\

Catálogo nuevo de representaciones & $C_4$ & Componente nuevo de VI4 & Par heterogéneo 2D--homogéneo equivalente; conservador; optimista. & Define la regla de equivalencia y el par que se usará para calcular la discrepancia de ampacidad. \\

Expediente de comparación & VI4 & Variable independiente discreta de OE4: $\{\mathrm{VD3},C_4\}$ & Combinaciones válidas de una VD3 concluyente con una realización de $C_4$. & Constituye la entrada completa de OE4 y produce VD4. \\

Contraste térmico & $C_k$ & Parámetro físico de caso; no VI ni VD & Bajo; medio; alto, con magnitudes fijadas en el piloto para cada realización de $E_3$. & Cubre diferencias de conductividad entre la región heterogénea y el entorno base y permite interpretar los indicadores contenidos en VD3 y VD4. \\

Extensión de la región heterogénea & $f_h$ & Parámetro físico de caso; no VI ni VD & Pequeña; media; grande, expresada como fracción del entorno local o tamaño normalizado. & Cubre el tamaño de la barrera o mejora térmica dentro de cada realización de $E_3$. \\

Proximidad al conductor & $d_h/D$ & Parámetro físico de caso; no VI ni VD & Cercana; media; lejana, normalizada por el diámetro exterior del cable o una distancia característica $D$. & Cubre la ubicación de la región respecto de la zona de mayor flujo térmico. \\

Condiciones pareadas & --- & Controles; no VI & Geometría, pérdidas, contornos, dominio, $T_{\mathrm{amb}}$ y configuración numérica aceptada. & Se mantienen constantes dentro de cada comparación para interpretar el efecto del atributo físico o de la representación. \\

Semilla o repetición numérica & $S$ & Metadato de robustez; no VI & Repeticiones seleccionadas en el piloto y en casos extremos. & Cuantifica variabilidad numérica de las salidas sin convertir el entrenamiento en una variable del estudio. \\

\end{longtable}

\tesisbodyfont

La malla máxima de cobertura combinará tres niveles de $C_k$, $f_h$ y $d_h/D$ dentro de cada una de las cuatro realizaciones heterogéneas simples de $E_3$. Esto produce hasta $3^3=27$ configuraciones paramétricas por realización.

Si se recorren completas esas cuatro realizaciones, el máximo de planificación será de 108 escenarios heterogéneos. Cada VD3 concluyente se combinará con la realización de $C_4$ que corresponda para formar VI4 y construir su comparación homogénea. El número 108 describe cobertura paramétrica potencial; no es el número de estados de VI3 o VI4 ni implica inferencia estadística factorial.

Este número se entiende como marco de planificación. Por tanto, si la carga computacional del piloto supera los recursos disponibles, se conservarán completos los patrones más críticos y se justificará cualquier fracción usada.

La lectura de resultados se hará en tres capas. Primero, se compararán los indicadores contenidos en VD3 y VD4 entre realizaciones discretas de $E_3$ y $C_4$, y se asignará a cada expediente uno de sus estados previstos.

Segundo, se describirá la sensibilidad física dentro de cada realización de $E_3$, especialmente para las combinaciones $C_k\times d_h/D$ y $C_k\times f_h$. Estas expresiones identifican cruces de parámetros del caso, no interacciones entre variables independientes metodológicas.

Tercero, se contrastará cada escenario heterogéneo con su homogéneo equivalente. Así se identificará cuándo la simplificación introduce una diferencia térmica u operativa relevante.

Para mantener consistencia con la matriz de operacionalización, las comparaciones pareadas mantendrán fijos geometría, pérdidas, contornos, dominio y configuración numérica. Por tanto, el efecto observado se atribuirá al atributo físico o al tipo de representación que diferencia el par.

Las condiciones no controladas se registrarán como metadatos. Entre ellas estarán la semilla, la tolerancia FEM publicada y la incertidumbre de las propiedades térmicas.

